\section*{Chapter \ref{ch:nw:fixed-income-and-rfq-platform-connectivity} --- Fixed Income and RFQ Platform Connectivity}

\begin{solution}{exo:nw:fixed-income-and-rfq-platform-connectivity:1}
It is not shown to the client: the dealer has no price in that auction, whatever it would have been.
\end{solution}

\begin{solution}{exo:nw:fixed-income-and-rfq-platform-connectivity:2}
U.S. markets in CME's NY5 in Secaucus, recovering in Aurora; European markets in LD4.2 in Slough, recovering in NY5.
\end{solution}

\begin{solution}{exo:nw:fixed-income-and-rfq-platform-connectivity:3}
The automatic stages sum to tens of milliseconds and practically never reach two seconds; a person, at six seconds' median, misses
two seconds 98.6\% of the time, so the chance is about $0.02 \times 0.986 = 1.97\%$, 1.94\% in the simulation.
\end{solution}

\begin{solution}{exo:nw:fixed-income-and-rfq-platform-connectivity:4}
At expiry the client compares every answer received in time whatever its order; the dealer's answers already arrive in time except
the person's, which a faster pricer does not change.
\end{solution}

\begin{solution}{exo:nw:fixed-income-and-rfq-platform-connectivity:5}
With five dealers asked, the best of the first five is the best of all who answer in time, unless someone answers after the timer: it is
``best at expiry'', except that execution happens as soon as the fifth answer arrives. With more dealers asked, it becomes ``best of the
first five''.
\end{solution}

\begin{solution}{exo:nw:fixed-income-and-rfq-platform-connectivity:6}
Treasury futures match in Aurora; with a cash book there too, both legs of a basis trade are in one building, and the New Jersey--Chicago
corridor's milliseconds and its costs drop out of the trade.
\end{solution}

\begin{solution}{exo:nw:fixed-income-and-rfq-platform-connectivity:7}
At about \qty{50}{\milli\second}, the competitors' own pricing time: 19.0\% under ``first acceptable'' against 19.2\% at expiry; faster
than the competitors, ``first acceptable'' favours the dealer, slower, it penalises it.
\end{solution}

\begin{solution}{exo:nw:fixed-income-and-rfq-platform-connectivity:8}
The timer is not the race when clients execute the first acceptable or the best of the first few answers: then the dealer races its
competitors, whose pricers are tens of milliseconds, and ten milliseconds are worth nearly fifteen cents of price.
\end{solution}

\begin{solution}{pb:nw:fixed-income-and-rfq-platform-connectivity:1}
\textbf{Part I.}
\begin{enumerate}
\item 1.94\%, almost all from the 2\% of requests answered by a person.
\item Zero at two seconds; 2.4\% at a \qty{100}{\milli\second} timer.
\item Above about half a second: the simulation finds 0.002\% at \qty{562}{\milli\second} and none from one second.
\item Both are no trade, but a missed answer also tells the client the dealer did not answer: a count against it in the client's dealer
  rankings.
\end{enumerate}
\textbf{Part II.}
\begin{enumerate}[resume]
\item 19.2\% at expiry, 26.4\% best of the first three, 39.0\% first acceptable.
\item 0, 1.6\% and 5.7\% of requests.
\item At expiry every answer received in time is compared, so the best acceptable one is always chosen.
\item It is faster than its four rivals (20 against \qty{50}{\milli\second} of pricing), so it is often the first acceptable answer.
\end{enumerate}
\textbf{Part III.}
\begin{enumerate}[resume]
\item 0.0, 2.3 and 12.6 points.
\item 0.9, 1.0 and 0.9 points.
\item About fifteen (12.6 against 0.9 points a cent).
\item From its own history: wins and losses by its response time and its rank in arrival order, per client; a win rate that falls with
  rank at equal price reveals a rule other than expiry.
\end{enumerate}
\textbf{Part IV.}
\begin{enumerate}[resume]
\item \emph{Named result}: with a pipeline whose automatic stages take tens of milliseconds and 2\% of requests answered by a person,
  the dealer misses 1.94\% of two-second timers; under a first-acceptable rule it loses 5.7\% of all requests to faster dealers despite
  the best acceptable price, and halving its pricing time gains 12.6 points of win rate against 0.9 for a cent of price.
\item Published: the platforms' protocols and automation, the order books' sites. Assumed: every pipeline number, the person's share and
  times, the timer, the price model's parameters, and the rules' use.
\item On the pricing stage if its clients use first-acceptable or first-$k$ rules; on the share sent to a person if timers are short; on
  price if they wait for expiry.
\item The advantage disappears and the win rates return toward one in five; speed then protects share rather than winning it.
\item Reduce it by widening the auto-responder's limits where the risk allows, and route the rest to a person with an alert and a
  deadline well inside the timer.
\item Everything: an order book is an exchange, with feeds, gateways and \omterm{def:nw:colocation-products-and-how-they-are-sold:colo}{colocation}, and microseconds matter again.
\item In the fixed income rows: platforms and sessions, the auto-responder's latency budget, the person's queue, and the order books'
  sites and \omterm{def:nw:colocation-products-and-how-they-are-sold:mmr}{cross-connects}.
\item When the client's rule executes before the timer, on the first acceptable or the first few answers.
\end{enumerate}
\end{solution}

\begin{solution}{iq:nw:fixed-income-and-rfq-platform-connectivity:1}
Network in, FIX engine, request enrichment, pricing (fair value, curve, comparables, inventory), limits and risk, answer out. Pricing,
if it recomputes a curve or waits for data; a person, if the request is outside the automatic limits.
\iqlookfor{Stages; the pricer's variance; the manual path.}
\end{solution}

\begin{solution}{iq:nw:fixed-income-and-rfq-platform-connectivity:2}
To execute quickly while the market is still where it saw it, to reduce information leakage from a long auction, and because its own
rules or workload call for automatic execution once the price is good enough.
\iqlookfor{Speed of execution; leakage; automation; the price tolerance.}
\end{solution}

\begin{solution}{iq:nw:fixed-income-and-rfq-platform-connectivity:3}
Log each request's arrival, the answer time, the rank in arrival order if the platform tells you, the price against the winning or cover
price, and the outcome; model win probability on price and response time together, per client.
\iqlookfor{Joint model of price and time; per-client rules; the cover price.}
\end{solution}

\begin{solution}{iq:nw:fixed-income-and-rfq-platform-connectivity:4}
Per-stage latency to find the stage; the pricer's data dependencies (a slow curve rebuild, a cache miss), garbage collection or
allocation, logging, the FIX engine's threads, and whether the share sent to a person changed.
\iqlookfor{Stage attribution; data dependencies; runtime effects; the manual share.}
\end{solution}

\begin{solution}{iq:nw:fixed-income-and-rfq-platform-connectivity:5}
It depends on the clients' execution rules: estimate the marginal win rate of a cent of price and of a millisecond from the dealer's own
history, and weigh them against their costs: a cent is given away on every trade, a millisecond is paid for once.
\iqlookfor{Rules decide; marginal values; recurring against one-off costs.}
\end{solution}

\begin{solution}{iq:nw:fixed-income-and-rfq-platform-connectivity:6}
FIX sessions to each platform with redundant lines; one pricing service with a latency budget and a fast path for liquid bonds; limits
and risk in-line; a person's queue with alerts; per-platform monitoring of misses, ranks and losses to speed; \omterm{def:nw:colocation-products-and-how-they-are-sold:colo}{colocation} and
\omterm{def:nw:colocation-products-and-how-they-are-sold:mmr}{cross-connects} for the order book in its data centre, with its own feed handlers and gateway.
\iqlookfor{Sessions; pipeline budget; manual path; monitoring; order book as an exchange.}
\end{solution}
